\section{模型评价与推广}

\subsection{模型优点}

\begin{enumerate}
	\item \textbf{分层预测兼顾异质性与一致性：}
	问题一由三级预测恢复18类交叉需求；验证期、测试期WAPE为0.7241、0.6519，84个滚动窗口均优于24小时同刻基线。
	
	\item \textbf{多目标调度能够刻画实际权衡：}
	问题二采用固定定标，50000个任务经全时域复算满足硬约束；成本降低507.34万元、新能源利用率提高0.223个百分点，同时保留碳排放和时延增加的代价。
	
	\item \textbf{储能模型保持题设负荷口径：}
	问题三固定题设负荷，仅优化能源侧；相较无储能方案，成本降低3936.19万元，碳排放、峰值和波动均下降，SOC及终端约束通过复核。
	
	\item \textbf{联合调度体现算储电协同效应：}
	问题四将顺序任务方案和联合任务方案按统一能源口径复算。联合任务方案使成本降低65.67万元、平均网络时延降低2.9851 ms、新能源未利用率降低0.0258个百分点；碳排放和峰值分别增加41.08 tCO$_2$、30.61 MW，完整呈现了服务质量改善的能源侧代价。两套任务方案各67项独立复核均通过。
\end{enumerate}


\subsection{模型不足}

\begin{enumerate}
	\item \textbf{预测依赖近期需求结构稳定：}
	18类需求恢复依赖历史结构，且相较直接预测的误差变化未达统计显著（$p=0.2163$）；业务突变时需采用动态窗口或时变权重。
	
	\item \textbf{问题二滚动启发式不能保证全局最优：}
	困难窗口与精确模型的最大标准化目标偏差为293.40\%；可按资源紧张度扩大邻域并延长关键窗口求解时间。
	
	\item \textbf{联合模型对极端情景和不确定性考虑有限：}
	问题三、四采用确定性电价、碳强度和新能源序列，未计预测误差与储能衰减；可引入逐级收紧约束及鲁棒或随机优化。
\end{enumerate}


\subsection{推广}

该框架可用于多区域数据中心、算力网络和园区源--荷--储调度，但须按实际任务时窗、网络、容量、PUE、电价、碳强度、新能源、储能和SLA重新标定。若引入任务抢占、电力互济或预测误差，还需扩展状态与不确定性约束并重新检验；可推广的是建模链路，而非本文参数和数值。
